\documentclass[12pt]{article}

% ------------------------------------------------------------
% Packages
% ------------------------------------------------------------
\usepackage[a4paper, top=2.5cm, bottom=2.5cm, left=3cm, right=2.5cm]{geometry}
\usepackage{helvet}
\renewcommand{\familydefault}{\sfdefault}
\usepackage{microtype}
\usepackage{parskip}
\usepackage{hyperref}
\hypersetup{
  colorlinks=true,
  urlcolor=blue,
  linkcolor=black
}
\usepackage{booktabs}
\usepackage{array}
\usepackage{xcolor}
\definecolor{accentblue}{RGB}{0,70,127}

% ------------------------------------------------------------
% Header rule
% ------------------------------------------------------------
\usepackage{fancyhdr}
\pagestyle{fancy}
\fancyhf{}
\renewcommand{\headrulewidth}{0pt}
\fancyfoot[C]{\small\thepage}

\begin{document}

% ============================================================
% Title block
% ============================================================
{\color{accentblue}\rule{\linewidth}{1.5pt}}
\vspace{4pt}

{\large\bfseries\color{accentblue} Cover Letter}

\vspace{2pt}
{\large Broadband Price Elasticity in EU and Eastern Partnership Countries:
Evidence of a Decade-Long Transformation (2010--2024)}

{\color{accentblue}\rule{\linewidth}{0.8pt}}

\vspace{8pt}
\noindent March 28, 2026

\vspace{12pt}
\noindent \textbf{Editor-in-Chief}\\
\textit{Information Economics and Policy}

\vspace{16pt}

% ============================================================
% Body
% ============================================================
Dear Editor,

We are pleased to submit the above manuscript for consideration as a research
article in \textit{Information Economics and Policy}. The paper addresses a
question central to both telecommunications economics and digital policy: has
broadband internet transitioned from a price-sensitive discretionary service to
a necessity good immune to price changes, and if so, when and where did this
transition occur?

\vspace{8pt}
\noindent\textbf{Contribution and scope fit.}\quad
Using a balanced panel of 33 countries --- 27 EU member states and 6 Eastern
Partnership countries --- over 2010--2024, we document a fundamental structural
transformation in broadband price elasticity. Eastern Partnership (EaP)
countries exhibited strongly elastic demand prior to 2020 (\(\varepsilon =
{-}0.60\), \(p < 0.001\)), while EU countries showed moderate elasticity
(\(\varepsilon = {-}0.10\)). By 2020--2024, both regions converged to
near-zero elasticity. Crucially, placebo tests reveal this transformation began
around 2015, predating COVID-19 and indicating a decade-long secular process of
digital integration rather than a pandemic-induced shock. This finding directly
advances the journal's mission of policy-oriented empirical research on
telecommunications markets, and extends a strand of work previously published
in this journal on broadband demand (Grzybowski et al., 2014; Liu et al., 2018;
Lindlacher, 2021; Sinclair, 2023) from the intensive margin (speed/quality
choices) to the extensive margin (the adoption decision itself) --- a dimension
for which time-varying evidence has been absent.

\vspace{8pt}
\noindent\textbf{Methodological contributions.}\quad
The paper makes two methodological contributions beyond the substantive
findings. First, we demonstrate that price measurement critically affects
inference in cross-country technology demand studies: income-relative prices (as
\% of GNI per capita) yield statistically significant EaP elasticities across
all eight control specifications, whereas PPP-adjusted prices --- the
conventional choice --- yield significance in only 12\% of EaP specifications.
This has implications extending beyond broadband research. Second, we employ
Driscoll--Kraay standard errors robust to heteroskedasticity, serial
correlation, and cross-sectional dependence, with 2SLS robustness checks
(first-stage \(F\)-statistic of 65--99) confirming that OLS estimates are not
contaminated by simultaneity bias.

\vspace{8pt}
\noindent\textbf{Policy relevance.}\quad
The EaP grouping is institutionally determined by the EU's Eastern Partnership
framework and the EU4Digital initiative, which targets exactly these six
countries for telecom regulatory harmonization with the Digital Single Market.
Our paper provides the first empirical test of whether demand-side convergence
between EaP and EU markets has occurred --- a prerequisite for assessing whether
full policy harmonization is feasible. The finding of converging near-zero
elasticity by 2020--2024 suggests that policy emphasis should now shift from
affordability subsidies to universal infrastructure deployment, with direct
implications for the EU Digital Decade 2030 targets.

\vspace{8pt}
\noindent\textbf{Replication.}\quad
All data are drawn from publicly available ITU and World Bank sources. Full
replication code and processed datasets are available in our GitHub repository,
as detailed in the Data Availability Statement.

\vspace{10pt}
This manuscript has not been published previously and is not under consideration
at any other journal. All authors have read and approved the submission. We
declare no competing financial interests or personal relationships that could
have influenced this work.

We believe this paper will be of strong interest to \textit{Information
Economics and Policy} readers and fits squarely within the journal's scope. We
are grateful for your consideration and look forward to the review process.

\vspace{16pt}
\noindent Yours sincerely,

% ============================================================
% Author block
% ============================================================
\vspace{12pt}
\begin{tabular}{@{}p{0.95\linewidth}}
  \textbf{Samir Orujov} \textit{(Corresponding Author)}\\[2pt]
  Head of Statistics Division,
  Information Communication Technologies Agency (ICTA), Baku, Azerbaijan\\
  School of Business, ADA University, Baku, Azerbaijan\\
  Azerbaijan State University of Economics (UNEC), Baku, Azerbaijan\\
  Universit\'{e} de Bretagne Sud, Vannes, France\\
  \href{mailto:sorujov@ada.edu.az}{sorujov@ada.edu.az}\\[8pt]

  \textbf{Ilgar Ismayilov}\\[2pt]
  Head of Economic Analysis Division,
  Information Communication Technologies Agency (ICTA), Baku, Azerbaijan\\
  Faculty of Social Sciences, Charles University, Prague, Czech Republic\\[8pt]

  \textbf{Jeyhun Huseynzade}\\[2pt]
  Deputy Chairman,
  Information Communication Technologies Agency (ICTA), Baku, Azerbaijan\\
\end{tabular}

\vspace{6pt}
{\color{accentblue}\rule{\linewidth}{0.8pt}}

\end{document}
